\begin{solapartat}
\punts{1,0} La reacció ha de complir la conservació del nombre màssic nuclear i, a més, ha de conservar la càrrega. Amb aquestes premisses, la reacció de desintegració és la següent:
\[ {}^{131}_{53}I \longrightarrow {}^{131}_{54}Xe + {}^{x}_{y}? + ? \]

Es compleix la conservació del nombre màssic:
\[ 131 = 131 + x \quad\Rightarrow\quad x = 0 \]
I per a la conservació de càrrega:
\[ 53 = 54 + y \quad\Rightarrow\quad y = -1 \]

Per tant, la partícula que busquem és un electró (${}^{\;\,0}_{-1}e$). A banda d'aquestes reaccions, per complir la conservació d'energia emeten neutrins, que, per conservar el nombre leptònic, hauria de tractar-se d'una antipartícula. Per tant, serà un antineutrí de la família electrònica. La reacció completa és:
\[ {}^{131}_{53}I \longrightarrow {}^{131}_{54}Xe + {}^{\;\,0}_{-1}e + \bar{\nu}_e \]

\destaca{Alternativament} es pot escriure ${}^{\;\,0}_{-1}\beta$ o $e^-$ enlloc de ${}^{\;\,0}_{-1}e$.

Si no s'indica l'antineutrí, cal restar 0,25~p.

No cal que l'estudiant detalli les operacions relacionades amb la conservació de la càrrega i del nombre màssic. Són operacions molt senzilles que es poden fer de cap; per tant, només cal valorar si ha estat capaç de deduir correctament els diferents termes de la reacció.

\punts{0,25} Es tracta d'una desintegració $\beta^-$.
\end{solapartat}

\begin{solapartat}
\punts{0,5} L'activitat és el nombre de nuclis que es desintegren per unitat de temps:
\[ A(t) = -\frac{dN}{dt} = \lambda N \]
Necessitem el coeficient de desintegració $\lambda$ que obtenim del temps de semidesintegració:

$A(t) = \frac{A_0}{2} = A_0\,e^{-\lambda t_{1/2}}$ i per tant, $t_{1/2} = \frac{\ln 2}{\lambda}$.

I $t_{1/2} = 8,02\ \text{dies} = 6,93\cdot 10^{5}\ \mathrm{s}$.

Directament, obtenim $\lambda = \frac{\ln 2}{t_{1/2}} = \frac{\ln 2}{6,93\cdot 10^{5}\ \mathrm{s}} = 1,00\cdot 10^{-6}\ \mathrm{s^{-1}}$.

Així, el nombre de nuclis és:
\[ N = \frac{A(t)}{\lambda} = \frac{40\cdot 10^{6}}{1,00\cdot 10^{-6}} = 4,00\cdot 10^{13}\ \text{nuclis} \]
I la quantitat de I-131:
\[ m = 4,00\cdot 10^{13}\ \text{nuclis}\,\frac{1\ \text{mol}}{6,022\times 10^{23}\ \text{nuclis}}\,\frac{131\ \mathrm{g}}{1\ \text{mol}} = 8,70\cdot 10^{-9}\ \mathrm{g} \]

\punts{0,25} L'activitat al cap d'una setmana ($t = 7\ \text{dies} = 6,05\cdot 10^{5}\ \mathrm{s}$) és:
\[ A(t) = A_0\,e^{-\lambda t} = 40\cdot 10^{6}\,e^{-1,00\cdot 10^{-6}\cdot 6,05\cdot 10^{5}} = 2,18\cdot 10^{7}\ \mathrm{Bq} \]

\punts{0,5} Representació gràfica correctament escalada, amb títols i unitats:
\figura[0.55]{sol_fig1.png}
\end{solapartat}
